\chapter{Bibliotheca mundi, Hic sunt dracones}

\begin{quote}
    \textit{L'univers (que d'autres appellent la Bibliothèque) se compose d'un nombre indéfini, et peut-être infini, de galeries hexagonales, la distribution des galeries est invariable. Vingt longues étagères, à raison de cinq par côté, couvrent tous les murs moins deux. Chaque étagère comprend trente-deux livres, tous du même format. La Bibliothèque est totale, et ses étagères consignent toutes les combinaisons possibles des vingt et quelques symboles orthographiques (nombre, quoique très vaste, non infini), c'est-à-dire tout ce qu'il est possible d'exprimer, dans toutes les langues.}

    Jorge Luis Borges, \textit{La Bibliothèque de Babel}, dans \og Fictions \fg, 1944
\end{quote}

Dans nos deux chapitres précédents, nous nous sommes d'abord interrogés sur le début des carrières pirates, puis sur le lien entre le vécu des individus et la création de communautés patrimoniales pirates. Dans ce troisième chapitre, nous édifierons une typologie d'étude des actions de ces communautés, en étudiant les dépôts \footnote{Pour joindre les conceptes d'archives, de bibliothèques et de tous les espaces communautaires ne tombant pas exactement dans ces définition, Abigail de Kosnik utilise le terme de répertoire. En réalité, son utilisation du terme répertoire est plus vaste que ce que étudions ici (il s'agit d'un répertoire culturel), mais je pense que l'utilisation du mot dépôt permet de bien comprendre la nature de ces tiers-lieux clandestins.} que créent et utilisent ces communautés. La manière dont est conçue une plateforme de communication et ses fonctionnalités peut grandement influencer la manière dont les individus échangent l'information. Par exemple, dans les archives des skyblogs de la BnF, nous retrouvons plusieurs blogs dont le but est d’offrir un accès à des fichiers et des logiciels, mais il est difficile d'observer une dynamique communautaire concernant la piraterie culturelle, ses conséquences politiques et ses mises en place infrastructurelles. Sur les plateformes de blogging et de microblogging comme Skyblog et Twitter/X, les dynamiques communautaires sont atténuées par le fonctionnement individualiste de partage des plateformes. Pour réussir à être visible sur Twitter, Facebook, Skyblog, ou la plupart des réseaux sociaux modernes, il faut avoir une base de followers, des gens qui se sentent touchés et investis par les messages, les blagues et les informations qu'un.e utilisateurice particulier fait circuler. Reddit et Discord sont quant à eux des réseaux purement communautaires, où les discours d'une personne peuvent toucher toute une communauté sans avoir besoin de construire une notoriété préalable. Nous étudierons ici les différents modes de partage, d'organisation et de maintien des dépôts pirates de différentes communautés.

\section{Naviguer sur les sept mers : une visite expresse des dépôts pirates personnels.}

Pourquoi ne pas utiliser le terme de bibliothèque pirate comme le fait l'article \textit{Lessons from the Library: Extreme Minimalist Scaling at Pirate Ebook Platforms} de Martin Paul Eve ? Ou pourquoi ne pas utiliser le terme de \textit{Rogue Archives} comme l'utilise Abigail de Kosnik dans son livre éponyme ? La raison de l'utilisation d'un nouveau terme à définir apparaît à la nature d'un constat, les archives des un.es peuvent devenir les bibliothèques et les répertoires des autres. Nous nous concentrerons ici sur la cartographie du partage interpersonnel, en rencontrant des individus qui partagent leurs collections de biens culturels au gré de leurs envies. S'ils peuvent faire partie d'une communauté patrimoniale, la manière dont ils mettent en place leurs systèmes de partage de l'information reste un élément personnel. 

\subsection*{Le dépôt personnel}

Le dépôt pirate personnel est, comme son nom l'indique, créé, alimenté et maintenu en ligne par une seule personne. Cette personne peut choisir de maintenir ce dépôt pour son profit propre, ou de le partager ; comme votre bibliothèque de salon, le dépôt personnel est défini par l'unicité de son origine et non par le nombre de personnes qui y ont accès. Ces dépôts peuvent être un simple dossier où l'individu stocke ses livres ou ses jeux piratés, ou des dispositifs de stockage beaucoup plus complexes et onéreux.

\begin{quote}
    T - \textit{[Avec mon frère] au début, on s'est échangé des clés USB ou des disques durs pour se partager des films qu'on avait envie de voir du registre de l'autre. Mais c'était devenu assez fastidieux, et surtout on c'est éloigné géographiquement pour avoir chacun sa vie pro dans différentes villes. [...] Donc il s'est posé la question de pouvoir avoir accès au registre de l'autre et de tout mettre en commun. On s'est renseigné, c'était une possibilité pour nous d'avoir notre propre serveur en fait. Récemment on a lancé ce projet à fond. On est encore en train de faire des testes et tout, mais on a ce qu'on appelle un NAS \footnote{network attached storage}. Ce serveur c'est mon ancien pc qui a été dépiécé et un peu agencé pour une nouvelle utilisation. On a mis de quoi rajouter plein de disques durs avec des énormes capacités, donc en l’occurrence on a en totalité ... La capacité totale, on a 4 disques de 8 Tera, ça fait 32 \footnote{32 000 gigaoctets}. [...] Pour nous ça fait un peu court, parce qu'on savait que nos deux bibliothèques, nos deux registres, on va dire cumulés, on était déjà à plus de 4 ou 5 Tera, et là actuellement, le dernier chiffre aujourd'hui on est presque à 7 tera au total.}
\end{quote}

Lors d'évènements particuliers, le dépôt personnel peut être ouvert par son propriétaire pour en faire profiter son cercle, ou toute une communauté. Dans le cadre de Théo et de son frère, le dépôt créé contient tellement de films et la logistique matérielle et logicielle le permettant, qu'ils décident de l'ouvrir aux membres de la famille et à leurs ami.es.

\begin{quote}
    T - \textit{Nous, TrueNAS, c'est un choix de facilité car c'est gratuit, et au dessus on a mis Jellyfin qui nous sert de lecteur média depuis ce système-là. [...] On a regardé la carte graphique de mon ancien pc, c'est une GTX 950. Donc ça date, ce n'est pas vraiment très puissant, mais il se trouve que c'est largement suffisant pour faire du décodage de 5 lecteurs en même temps. [...] On ouvrirait à des amis et on leur ferait payer à l'année une petite somme, quelque chose comme 10 euros, simplement pour compenser la facture d'électricité d'avoir un pc qui tourne tout le temps. La deuxième motivation c'est pour nos parents, nos parentes étaient très dépendants, ils avaient opté pour la facilité, c'est que te propose par exemple Netflix. Pour eux c'est la même chose, la différence, c'est qu'ils ne paieront plus Netflix, et nous derrière on aura plus de contrôle sur ce qui se passe sur le catalogue pour éviter que des trucs disparaissent du jour au lendemain.}
\end{quote}

D'autres exemples moins technophiles peuvent être intéressants à étudier. Nous avions évoqué les vagues de partage de films de David Lynch. À l'annonce de la mort du réalisateur le 16 janvier 2025, de très nombreuses personnes se sont mises à télécharger, riper et partager les œuvres du réalisateur. En ouvrant son dépôt personnel en le transférant sur son Google Drive ou un MegaUpload, les individus prenaient part au deuil international de la communauté des fans. Le piratage et son partage revêtent ici à la fois une dimension mémorielle (les œuvres ne disparaissent pas si les fans en créent des copies) et communautaire (je peux contribuer à rendre hommage à David Lynch en propageant son art).

\begin{quote}
\textit{    coucou j’ai téléchargé toutes les œuvres de Lynch (films et série) et je les ai mis dans un mega pour que ce soit plus facile pour toutes les personnes qui veulent (re)découvrir ce qu’il a fait :) et n’oubliez pas, si vous commencez twin peaks l’ordre c’est : saison 1 —> saison 2 —> Fire walk with me —> saison 3
}

    @cycylcds, sur twitter/X, le 18/01/2025
\end{quote}

\subsection*{Les individus \og passerelles \fg et la culture archivistique amatrice}

Il existe un entre-deux entre le dépôt personnel et le dépôt communautaire. Il s'agit de dépôts dont l'administrateur est une personne particulière, mais dont le but de création est d'être partagé à une échelle communautaire. Ces dépôts sont parfois éditorialisés, par des individus qui sont reconnus dans leurs communautés comme étant des ressources précieuses de savoir et de connaissance. Ce genre d'individu que nous pouvons surnommer des \og passerelles \fg sont relativement courant dans les communautés politisées comme les milieux militants et LGBT. Ce partage passe par des infrastructures et une organisation plus pérennes que les simples dépôts personnels. Nous pouvons trouver dans cette catégorie des exemples divers, allant de la subversion de systèmes électroniques pour la réparation d'objets, à la mouvance du \textit{bio-hacking} et de l'Anarcho-transhumanisme transgenre refusant d'être dépendant.es de l'ordre médical et de l'industrie pharmaceutique pour leurs transitions de genre. 

\begin{quote}
    \textit{They're both on my website now!! (Une image représentant deux fascicules) Transfem DIY HRT, Transmasc DIY HRT, Nobody can stop you}

    @desert\_varnish, sur twitter/X sur 02/02/2025
\end{quote}

La principale communauté de personnes \og passerelles \fg que nous avons pu étudier se situe sur le forum subreddit r/DataHoarder. Nous l'avons déjà citée dans le chapitre précédent sans vraiment expliquer en profondeur par quoi cette communauté se passionne ainsi que comment interagissent ces membres. Le terme de datahoarder provient du mélange du terme data (données) et hoarder (ceux qui stockent). Cette communauté est composée de personnes fascinées par le stockage, l'archivage et la mise à disposition de jeux de données. Les datahoarders mettent en avant les qualités de pérennisation et de partage de leurs bases de données. Cependant, le choix des données qu'ils stockent et la mise à disposition est un choix personnel, qui suit souvent les intérêts et passions du hoarder. r/Datahoarder est donc un rassemblement de personnes stockant et partageant des collections de données rares, mais le fonctionnement interpersonnel (demander à un hoarder s'il compte rendre ses données publiques) fait que nous sommes face à une multitude de bibliothèques personnelles plutôt qu'à une communauté organisée autour d'un outil.

Il est passionnant de se plonger dans les interactions d'une communauté si riche qui se base sur l'interaction sur un forum ouvert pour communiquer. Par exemple, le forum r/DataHoarder permet aux membres de créer une étiquette qui s'affiche après le pseudonyme ; ils s'en saisissent afin d'afficher les différentes tailles de leurs collections et infrastructures. Il se dégage ainsi une forme de hiérarchie immanente au volume qu'une personne stocke et à la logistique et l'argent qu'iel met en place pour mener à bien sa mission de hoarder. Certaines étiquettes laissent deviner une forme de syndrome de Diogène numérique, avec des personnes indiquant avoir mis en place des serveurs de stockage de plusieurs centaines à plusieurs milliers de téraoctets de stockage. Si beaucoup de ces utilisateurs sont des ardents défenseurs du stockage et du piratage de contenus culturels, beaucoup ne se considèrent pas comme des pirates. La citation ci-dessous est particulièrement intéressante car u/Blue-Thunder utilise la terminologie académique du livre \textit{Rogue Archives} d'Abigail de Kosnik, en modifiant la définition qu'iel applique derrière \footnote{Agigail de Kosnik utilise ce terme de \textit{Rogue Archivist} pour parle principalement des personnes archivants et partageants des contenus créer par des communautés de fan (fanfictions, fanarts), ici u/Blue-Thunder archive des fichiers qui sont des dossiers sous droit, comme des films, des livres, des vidéos, ce qui est très différent du première exemple. }. 

Un excellent exemple trouvé par Marina Hervieu avec les outils de recherche de la Bnf est le skyblog \textit{brotjackta} de l'utilisateur Ruinessi. À première vue, ce skyblog est rempli de posts vides qui affichent la phrase suivante \textit{Sorry, the page you requested has been moved.}. Les titres des posts sont pourtant très intéressants, il s'agit de tutoriels anglophones, pour découvrir plusieurs outils de rapatriement de données torrents ou de téléchargements. En réalité, le blog de Ruinessi est écrit en blanc sur fond blanc. Avec cette astuce, il espère éviter d'attirer l'attention des personnes n'ayant un minimum de connaissances en navigation web et informatique. Nous avons trouvé d'autres exemples d'individus assurant un rôle de passerelle dans les archives des skyblogs. Les posts et les blogs ont la particularité d'être très simplement personnalisables, avec notamment du HTML, et de pouvoir sauvegarder une petite quantité de données extérieures comme un fichier mp3 ou quelques images. Certains utilisateurs ont profité de ces possibilités, couplées à la capacité d'offrir des hyperliens vers des sites extérieurs pour offrir des mini-bibliothèques pirates.

Voici quelques exemples de blogs de bibliothèques pirates de la plateforme Skyblog.
\begin{itemize}
    \item \textit{FilmstreamingVF}, créé en 2016 par l'utilisateur telechargerfilm, probablement lié à un site de téléchargement filmvfr.com. Ce blog utilise des petites fiches d'informations techniques, ainsi que des critiques de différents films de l'actualité cinématographique pour amener ses utilisateurs à passer sur son site de streaming pirate.
    \item \textit{téléchargergratuitement}, créé en 2018 par un utilisateur s'inspirant du fonctionnement de 01.net,  qui est un peu particulier ; en effet, ce site propose de télécharger des applications de manière légale. Son modèle de fonctionnement s'inspire des applications d'installation comme les différents gestionnaires de paquets sous Linux. Chaque article de blog s'intitule selon une application, avec une description de son utilité et de son usage. 
    \item \textit{La biblio de Tetel} et \textit{Le monde de Tetel}, deux blogs liés et créés par l'utilisateur Kurisu/Tetelle dont la première capture sur la Wayback Machine remonte à 2013. Cette personne reprend les codes de l'article de blog en donnant un avis critique sur plusieurs œuvres, en ajoutant des liens de téléchargement dans le corps du texte ou dans les commentaires. 
    \item \textit{Mangas-a-telecharger}, créé en 2008 par un utilisateur dont je n'ai pas trouvé le pseudonyme. Ce blog est intéressant car il s'agit du blog de téléchargement pirate le plus assumé de nos trouvailles. Après une petite introduction, l'auteurice du blog affiche une liste de liens pour avoir accès à ses mangas favoris. Quelques posts de blog demandent de l'aide à ses abonnés pour trouver les mangas sur lesquels iel n'arrive pas à mettre la main. 
\end{itemize}

Nous pouvons également trouver quelques initiatives solitaires comme le blog \textit{Anti-téléchargement} \footnote{https://web.archive.org/web/20130517194237/http://anti-telechargement.skyrock.com/} essayant de prévenir et de sensibiliser aux dangers du piratage, ou le blog \textit{Unis-Contre-Le-Plagiat} \footnote{https://web.archive.org/web/20100215052912/http://unis-contre-le-plagiat.skyrock.com/} sur la question du plagiat de posts entre utilisateurs de Skybklog. \\

La plupart de ces blogs prennent souvent la même forme. Un post de blog avec comme titre le nom de l'œuvre (souvent des mangas ou des films) avec ensuite un résumé et une critique de l'œuvre, puis une note et enfin un lien de rapatriement du fichier.\\

%ajouter un élément qui introduit cette citation au mon dieu 

Sur r/DataHoarder, les habitués du forum se définissent eux-mêmes comme des archivistes rebelles, revendiquant un fonctionnement clandestin au nom d'une lutte pour le patrimoine numérique. 

\begin{quote}
    \textit{I'm sorry but I think this terminolgy is outdated. I prefer to be called a Rogue Archivist. I'm not on the high seas pludering loot from corporations, I am protecting media from vanishing forever, against the will of the respective copyright owners who would love to see certain things just vanish, or who are so incompentent that they do not have (redundant) backups.}

    u/Blue-Thunder, en réponse au poste de u/pilimi\_anna \og  How to become a pirate archivist \fg sur r/datahoarder, 2022
\end{quote}

Cependant, r/DataHoarder n'est pas dédié au stockage de données dans un but précis, seulement au stockage et à la mise à disposition de données de masse, et ceci peut amener à des discussions de questions éthiques. Si le stockage de contenus culturels piratés n'est pas considéré comme un problème ou un dilemme, le stockage et la mise à disposition de données personnelles l'est beaucoup plus. Par exemple, plusieurs projets de membres se concentrent sur le scrapping, le stockage et la diffusion de données pornographiques. Le \textit{Petabyte Porn Project} \footnote{u/-Archivist, \textit{The Petabyte Porn Problem: Public webcam social interaction archiving. [Chaturbate, MyFreeCams, Streamate]}, sur r/DataHoarder, 2017} se présente sous l'angle de l'étude et de l'archivage des \og interactions sociales \fg. Pour ce faire, les individus qui se joignent au projet via une feuille de contact peuvent archiver dans une base de données commune les fichiers vidéo de plusieurs sites de pornographie en direct \footnote{\og Le phénomène connu sous le nom de « caming », qui consiste pour des « camgirls » ou « camboys » à proposer, moyennant rémunération, une diffusion d'images ou de vidéos à contenu sexuel, le client pouvant donner à distance des instructions spécifiques sur la nature du comportement ou de l'acte sexuel à accomplir, n'entre pas dans le cadre de la définition de la prostitution qui implique un contact physique onéreux avec le client pour la satisfaction des besoins sexuels de celui-ci.\fg, citation d'un acte de la chambre criminelle de la cour de cassation, n° 21-82.283 B, 18 mai 2022} (Chaturbate, MyFreeCams, Streamate). Plusieurs utilisateurs vantent les utilisations possibles d'une telle collection : étude sociologique sur les interactions entre regardants et regardé·es \footnote{u/t0asterB0y \og This could be an invaluable resource for social scientists and behavioral psychologists. \fg, 2017}, tentative d'application de reconnaissance faciale pour retrouver des victimes de traite humaine \footnote{u/deleted,  \og Have you thought about doing facial recognition and matching up to sex trafficking photos, missing persons photos, etc? \fg, 2017}. Cependant, r/DataHoarder est une communauté essentiellement composée d'hommes, et les réponses à ce post de blog sont remplies de personnes demandant à avoir accès à la base de données pour chercher leurs stars préférées, en particulier celles qui ont arrêté de gagner leur vie avec la pornographie et ont supprimé leurs vidéos en espérant ne plus avoir à assumer une notoriété publique sur le web. D'autres éléments problématiques sont mis en avant pour les utilisateurs les plus sceptiques : la détention et le stockage de fichiers pornographiques de personnes mineures sont strictement interdits dans la grande majorité des pays, même à des fins documentaires et de tentative d'aider les autorités compétentes. 

\section{La carte et le coffre au trésor : les dépôts communautaires pirates}

Dans notre typologie des dépôts pirates, toutes les initiatives de stockage et de partage que nous avons étudiées ont comme points communs d'être des initiatives individuelles. Ces dépôts peuvent être personnels ou, au contraire, s'adresser à une communauté, mais les créateurs et les mainteneurs de ces projets étaient toujours des personnes seules, ou des petits groupes de personnes liés par des relations interpersonnelles. Cependant, ce n'est pas parce qu'un dépôt pirate est un dépôt personnel qu'il ne peut pas s'inscrire dans une communauté patrimoniale pirate. Ce qui différencie les dépôts personnels et les dépôts communautaires est à la fois le mode de fonctionnement, l'ampleur des projets et la cible de la diffusion des projets. 

\subsection*{Les dépôts communautaires pirates}

Les exemples les plus connus de dépôts communautaires sont les bibliothèques, académiques ou pas, pirates. Que ce soit les plus connues : Sci-Hub, LibGen et Zlibrary, le sujet de ces bibliothèques en ligne a été l'épicentre de nombreux débats, d'articles médiatiques et de quelques études. Les bibliothèques pirates comme LibGen sont des dépôts communautaires pour plusieurs raisons : 

\begin{itemize}
    \item Ces dépôts pirates sont trop massifs pour être gérées par un seul utilisateur, plusieurs équipes entières travaillent sur ces projets afin d'assurer la logistique, de maintenir le fonctionnement des sites et des serveurs, mais également de développer de nouvelles fonctionnalités pour les utilisateurs. Certaines de ces projets sont également open-sources, ce qui agrandit encore le cercle de personnes impliquées dans le fonctionnement. Prenons un exemple concernant le jeu-vidéo : le projet XenomRecomp du développeur Hedge-dev \footnote{https://github.com/hedge-dev/XenonRecomp} permet de recompiler des jeux Xbox 360 en des exécutables natifs pour ordinateurs Windows. Copier et recompiler du code est évidemment une pratique pirate, mais cet outil est aussi extrèmement efficace pour sauvegarder et faciliter l'utillisation de toute la librairie de la console. Les personne ayant des connaissances aboutis en développement, et de l'expérience avec la compilation ou l'émulation de logiciels peuvent aider le développeur original à fiabiliser et développer son outil grâce au système de tickets nommé \textit{Issues} intégré au site Github.

    \item Tous les dépôts dont nous avons étudiés dans cette partie n'ont pas été créé par et pour une seule communauté. Ils ont pour but d'être partager et utilisé par quiconque qui peut trouver une utilité dans leurs outils et les documents que les dépôts communautaires mettent à disposition. Plusieurs de nos exemples cités ont pour but de mettre à disposition le plus d'imprimés possibles et ce à l'échelle internationale. 
\end{itemize}

Zlibrary, comme SciHub, LibGen ou Anna's Archive, sont l'architecture de fonctionnement sur laquelle nous construisons le concept de \textit{l’organisation anarchique de guérilla} que nous avons évoqué dans notre introduction et notre deuxième chapitre. Ce sont des sites qu'il peut être difficile de trouver quand on ne sait pas où chercher. La plupart de ces exemples ne sont pas directement disponibles en France dû à des décisions de justice imposant aux principaux fournisseurs d'accès à Internet de les effacer de leurs DNS \footnote{
Domain Name Systemn ou Système hiérarchique de nommage des ressources d'un réseau : un des protocoles de base de l'Internet, sans inscription d'un site à un dns, vous ne pouvez pas vous y connecter.
}. Et certaines des URL de ces sites sont régulièrement fermées, affichant alors une image bien connue des habitués, une page bleue avec le logo du FBI américain ainsi qu'une explication de la saisie du domaine et du caractère illégal du site qui y était hébergé. Il faut alors pour l'utilisateur trouver des alternatives, les quatre principales bibliothèques ont ce que l'on nomme des miroirs, des sites jumeaux, ou des messageries automatisées sur le service Telegram, permettant de faire des demandes de retrait de fichiers sans passer par l'interface utilisateur du site. Ces solutions de repli peuvent prolonger la vie du site pendant quelques temps avant que les autorités ne remontent aux serveurs. Quand un service passe complètement hors ligne, nous pouvons observer des messages de déploration et de recherche d'alternatives sur les principaux forums communautaires (reddit) et les sites de microblogging (twitter/x).

\begin{quote}
    \textit{I miss Zlib everyday </3}

    @BLincoln56151, sur Twitter/X, le 20 avril 2025
\end{quote}

Cependant, les équipes travaillant sur ces bibliothèques d'Alexandrie contemporaine sont si vastes et organisées qu'une nouvelle version du site ne tarde pas à réapparaître en ligne après quelques semaines ou quelques mois. Cependant, il est difficile pour les équipes officielles de ne pas concurrencer des groupes malveillants de phishing, cherchant à récupérer des adresses e-mails et des mots de passe en créant des faux clones incomplets et en les faisant passer pour le site officiel.

\subsection*{Bibliothèques pirates et archives rebelles : étude de cas avec \textit{Zlibrary} et \textit{The Eye}}

\subsubsection*{\textit{Zlibrary}}
Prenons un exemple : la plateforme Zlibrary est une bibliothèque clandestine mettant à disposition des milliers de livres et d'articles sur tous les sujets. Zlibrary correspond parfaitement aux mécanismes décrits dans le chapitre précédent de simplicité apparente, avec un minimalisme de surface qui peut créer une impression d’inhospitalité à première vue. La première page du site se compose du logo de la Zlibrary, suivi du slogan \og \textit{ Your gateway to knowledge and culture. Accessible for everyone. } \fg et d'une simple barre de recherche. Mais si l'on passe cette première introduction minimaliste, Zlibrary offre énormément de fonctionnalités bien plus avancées qu'un simple catalogue de livres à télécharger. 

Premièrement, le site propose un système de compte permettant aux utilisateurices de garder une trace des livres qui les intéressent, de créer des listes de lecture, de se rappeler de ce qu'iels ont téléchargé, d'être informé·es de leur nombre de téléchargements restants, d'avoir accès à des recommandations de lecture personnalisées selon leurs intérêts. Zlibrary \footnote{https://z-lib.gl/, mais vous n'avez pas le droit d'y aller, ce serait illégal} est un outil avec une présentation minimaliste, mais qui essaie de fournir à ses utilisateurices une expérience lisse, complète, parfois plus adaptée à certains publics. Par exemple, pour le public estudiantin et les professionnels qui ont besoin d'un accès à la littérature académique, Zlibrary est un outil beaucoup plus adapté que n'importe quelle plateforme de vente d'ebook. On y trouve à la fois de grandes monographies, différents éditeurs, des revues, des recommandations.

Ensuite, on y découvre l'outil des listes de lecture créées par d'autres utilisateurices. Ces listes de lecture pourraient être un objet d'étude à part entière ; par défaut, les comptes Zlibrary sont publics et tous les utilisateurices peuvent voir ce que les autres utilisateurices regardent, classent et téléchargent. Imaginons un utilisateur type de Zlib, un étudiant travaillant sur le piratage comme vecteur d'une culture archivistique populaire. Pour garder une trace des différents livres et papiers qu'il explore sur ce sujet, il crée plusieurs listes de lecture, la première se nomme \textit{Piracy}, la seconde \textit{Digital Curation}, la troisième \textit{Digital Humanities}. Il utilise alors Zlibrary comme une bibliothèque, afin de mener à bien une recherche bibliographique. Une fois son travail de recherche bibliographique terminé, ses documents téléchargés, ses listes de lecture publiques perdurent, elles deviennent des archives numériques, des traces de son sérieux académique et de son érudition. Cependant, le caractère public et indexé dans le moteur de recherche des listes permet de leur donner une seconde vie. Plusieurs autres utilisateurices vont également faire des recherches sur ces sujets pour leurs études ou par curiosité, et alors les archives des uns deviennent les bibliothèques des autres. Au total, archives de lecture de notre étudiant exemple ont été consultées 586 fois. Si Zlibrary offre des outils de classification et de recherche avancée, les listes de lecture sont une fonctionnalité précieuse qui permet de se reposer sur le travail bibliographique d'autres personnes pour accéder à des documents moins populaires, moins cités, ou qui ne sont pas considérés comme pertinents par l'algorithme de recherche.\\

\subsubsection*{\textit{Archive Team \& le projet The Eye}}

Archive Team \footnote{https://wiki.archiveteam.org} est une communauté en ligne ressemblant beaucoup dans ses missions et son fonctionnement à r/DataHoarder. Le site d'Archive Team se présente comme une équipe de \og d'archivistes rebelles \fg, de développeurs et d'écrivains qui se mobilisent pour préserver l'histoire et l'héritage numérique depuis 2009. Cette communauté utilise et reformule la définition de \textit{d'archivistes rebelles} définie par Abigail de Kosnik. Comme r/DataHoarder, les projets d'Archive Team se sont construits sur la réflexion de l'enjeu culturel et politique de la conservation de certaines données par des groupes qui ne sont ni des institutions, ni des entreprises. Le projet \textit{The Eye} mis en place par l'Archive Team est cependant très différent d'un dépôt pirate en centralisant les actions de ses membres en un seul collectif. Si le projet d'Archive Team se présente comme étant un groupe respectant les règles établies de la propriété intellectuelle, l'un des projets issus de cette communauté, \textit{The Eye} s'est fait connaître pour son archivage et sa mise à disposition de dépôts de torrents pour l'entièreté de Library Genesis (LibGen) en 2019 et nous pouvons également retrouver dans l'arborescence \textit{../files/comics/..} des archives complètes de mangas sous droits, mais également des manuels de jeu de rôle, des articles scientifiques pour de l'entraînement DIY de modèles d'intelligence artificielle, et des collections comme \textit{The Telegram [Anti]Social Research Archives}\footnote{https://the-eye.eu/tasra/} avec plus de 200 groupes de discussions publics sur la plateforme Télégram sélectionnés pour la virulence de leurs membres avec presque 4 To de données réunies. Le site reprend énormément du fonctionnement et de l'identité graphique des dépôts pirates qui revendiquent leur appartenance à la mouvance des dépôts pirates. Le simple nom de la communauté peut sembler inquiétant : L'œil, avec un logo représentant un œil ouvert, omniscient, au centre d'un triangle équilatéral où chaque segment est marqué de trois marques. La page d'accueil est composée d'un fond noir, avec en transparence des paires d'yeux dessinés avec une direction artistique ésotérique et les quelques lignes de texte écrites ne permettent pas à l'utilisateur néophyte de continuer sur le site pour en apprendre plus sur l'organisation, le fonctionnement et les projets en cours de The Eye et de l'Archive Team. Cette présentation minimaliste récurrente, couplée à des directions artistiques souvent ésotériques, monochromes, chaotiques, encourage une expérience utilisateur de curiosité et de jeu avec un interdit réel ou fantasmé \footnote{Martin Paul Eve, \textit{warez: the infrastructure and aesthetics of piracy}, édition puctum books, 2021}. L'esthétique et le politique se mélangent ici, pouvant notamment nourrir la manifestation d'une représentation du monde que nous avions définie comme un zèle ou un complotisme archivistique dans le chapitre précédent. 

\section{Conclusion}

Dans ce chapitre, nous avons produit un panorama des différents types de dépôts pirates, en discriminant entre les différentes architectures et organisations des individus qui les créent et les maintiennent. Nous avons appliqué les grilles de lectures et d'analyses de Martin Paul Eve qui mélangent les modalités d'organisation de ces dépôts avec une suite de codes culturels, une esthétique du piratage. Ceci nous a permis de comprendre la traduction du principe d'organisation de guérilla et les influences sous-jacentes qui en font un référentiel d'organisation autant technique que politique. Ce référentiel est composé d'un discours composé de plusieurs sphères d'influences. Le cœur de ce discours est en grande majorité alimenté par la culture du libre et ses théoriciens/ingénieurs principaux comme Richard Stallman, Aaron Swartz ou Linus Torvalds. D'autres éléments peuvent être empruntés à la culture militante marxiste, anarchiste et libertaire. Comme toutes les cultures qui remettent en cause le statu quo, la culture technique et politique des individus baignant dans le monde des dépôts pirates peut amener à des tensions internes. Certains individus désirent conserver tout ce qui est à la portée de leurs capacités de stockage, y compris des données personnelles sensibles, ou en entrant dans une forme de \og complotisme archivistique \fg. Le fonctionnement, l'organisation et le rôle de ces dépôts pirates peuvent être interprétés comme le symptôme d'une société qui cherche à se défendre face à l'organisation du travail et du rôle de la marchandise dans son fonctionnement. Nous pouvons citer ici les grandes lignes des théories de la littérature marxiste sur le sujet de la culture, sa place dans la superstructure et son lien avec l'infrastructure. À l'époque de la reproductibilité technique parfaite de la grande majorité des objets culturels, ces dépôts pirates, les individus et les communautés qui les maintiennent, jouent un rôle important dans la sauvegarde et la diffusion de savoir et portent en eux la graine d'un possible d'une distribution plus juste de la culture et d'une collectivisation de l'information. 


%En conclusion, les dépôts pirates, qu'ils soient personnels ou communautaires, jouent un rôle essentiel dans la sauvegarde et la diffusion du savoir à l'ère numérique. Les dépôts communautaires, tels que LibGen, Sci-Hub et Zlibrary, démontrent une capacité impressionnante à s'organiser et à prospérer malgré les défis juridiques et techniques. Ces bibliothèques pirates représentent non seulement des efforts collectifs importants pour contourner les limitations imposées par les systèmes traditionnels d'accès à l'information, mais elles suscitent également des débats persistants sur la légitimité et l'éthique de leur existence. 

%En conclusion, la plateforme Zlibrary ainsi que des initiatives comme Archive Team et The Eye illustrent une approche unique et audacieuse de l'archivage numérique et de la diffusion du savoir. Zlibrary, avec son minimalisme apparent et ses fonctionnalités étendues, se présente comme une ressource précieuse pour ceux qui recherchent un accès universel à la connaissance. Simultanément, Archive Team et The Eye s'engagent dans la préservation et le partage des informations, tout en remettant en question les paradigmes traditionnels de la propriété intellectuelle et de l'archivage. Ensemble, ces plateformes témoignent d'une mouvance qui mêle esthétique, politique et une certaine subversion pour rendre accessible une vaste gamme de contenus. Cependant, cette accessibilité soulève des questions éthiques et légales qui méritent d'être examinées de près.


%%III/ Dépôts numériques, hybride entre un lieu de
%conservation et de partage.
%• Le dépôt pirate, entre archive et bibliothèques
%– La fonction de conservation des dépôts numériques.
%– Le partage et la diffusion des ressources piratées.
%• Les communautés autour des dépôts numériques
%– Les contributeurs et les utilisateurs des dépôts.
%¨ Les dynamiques sociales et les interactions au sein des communautés.